% !TeX root = ../main.tex
\chapter{Getting started}
\section{What is HobbySpark?}

HobbySpark is a programming language and development environment
designed for creating programs that interact with hardware such as
microcontroller boards and electronic components.

It provides a simple, readable programming language that lets you
write programs using familiar concepts such as variables, conditions,
loops, functions, and classes, while also providing classes designed
specifically for working with hardware.

Instead of requiring the user to deal directly with all of the
lower-level details of hardware programming, HobbySpark provides
higher-level objects and methods that represent the components being
used. This allows a program to describe what it wants the hardware to
do in a more natural way.

For example, a hardware component can be represented by an object:

\begin{lstlisting}
	led = Led(6)
\end{lstlisting}

The program can then interact with that object through its methods:

\begin{lstlisting}
	led.on()
	led.off()
\end{lstlisting}

HobbySpark is intended to make experimenting with electronics and
programming more approachable while still allowing users to build
increasingly complex projects.

\subsection{How HobbySpark Works}

From the user's point of view, the workflow is straightforward:

\begin{center}
	\textbf{Write your program}
	$\rightarrow$
	\textbf{Check it}
	$\rightarrow$
	\textbf{Build it}
	$\rightarrow$
	\textbf{Upload it}
	$\rightarrow$
	\textbf{Watch your hardware do something}
\end{center}

The user does not need to manually construct the underlying C++
program. HobbySpark handles the conversion required by its hardware
toolchain.

\subsection{What Can You Do with HobbySpark?}

HobbySpark programs can use ordinary programming features alongside
hardware-oriented functionality. This means the same language can be
used for both the logic of a program and the control of physical
hardware.

Some of the features available to HobbySpark programs include:

\begin{itemize}
	\item Variables
	\item \texttt{if}, \texttt{elif}, and \texttt{else} conditionals
	\item \texttt{while} and \texttt{for} loops(for loops only applicable for range)
	\item Functions
	\item Classes with constructors and operator overloading(with explicit return types)
	\item All python basic types (\texttt{bool}, \texttt{int}, \texttt{float} and \texttt{str})
	\item \texttt{break} and \texttt{continue}
	\item Programs that continuously interact with connected
	hardware
\end{itemize}

The goal of HobbySpark is simple: write code, connect hardware, and
make something happen.

\section{Requirements}
You will need the following:
\begin{enumerate}
	\item \textbf{Computer}
	\item \textbf{Applicable USB cable}
	\item \textbf{HobbySpark app}
	\item \textbf{Board}
	\item \textbf{The required devices like LEDs, motors, and sensors}
\end{enumerate}
